\exercisesection{\S~11}

\begin{enumerate}
\def\labelenumi{\arabic{enumi})}
\tightlist
\item
  令 \(\mathfrak{g} = \mathfrak{sl}(2, k)\)。置 \(G = \operatorname{Aut}_e(\mathfrak{g})\)；该群可与 \(\mathbf{PSL}_2(k)\) 相识别，参见第VII章§3第1节注2。
\end{enumerate}

\begin{enumerate}
\def\labelenumi{\alph{enumi})}
\item
  \(\mathfrak{g}\) 的每个幂零元素都 G-共轭于 \(\left( \begin{array}{cc}0 & \lambda \\ 0 & 0 \end{array} \right)\)（对某个 \(\lambda \in k\)）。这样的元素是主的当且仅当它非零。
\item
  元素 \(\begin{pmatrix}0&\lambda\\0&0\end{pmatrix},\begin{pmatrix}0&\mu\\0&0\end{pmatrix}\)（其中 \(\lambda,\mu\in k^{*}\)）G-共轭当且仅当 \(\lambda^{-1}\mu\) 是 k 中的平方。
\item
  \(\mathfrak{g}\) 的每个单元素都 G-共轭于 \(\left( \begin{array}{cc}1 & 0\\ 0 & -1 \end{array} \right)\)。
\end{enumerate}

\begin{enumerate}
\def\labelenumi{\arabic{enumi})}
\setcounter{enumi}{1}
\tightlist
\item
  令 \(A = \begin{pmatrix} 0 & 1 \\ 0 & 0 \end{pmatrix}\)，\(B = \begin{pmatrix} 0 & 0 \\ 1 & 0 \end{pmatrix}\)。则 \(A\) 是幂零的，而 \(AB\) 不是，且
\end{enumerate}

\[
[ A, [ A, B ] ] = \left( \begin{array}{c c} 0 & - 2 \\ 0 & 0 \end{array} \right).
\]

由此推出，引理5不能推广到特征 2 的域上。

\begin{enumerate}
\def\labelenumi{\arabic{enumi})}
\setcounter{enumi}{2}
\tightlist
\item
  令 \(\mathfrak{r}\) 为 \(\mathfrak{g}\) 的根基，\(\mathfrak{s} = \mathfrak{g} / \mathfrak{r}\)。证明下列条件等价：
\end{enumerate}

\begin{enumerate}
\def\labelenumi{(\roman{enumi})}
\item
  g 不含 \(sl_{2}\) 三元组；
\item
  \(\mathfrak{s}\) 不含 \(\mathfrak{sl}_2\) 三元组；
\item
  \(\mathfrak{s}\) 是各向异性的（§10习题6）；
\item
  \(\mathfrak{s}\) 不含非零的可对角化元素（§3习题10）。
\end{enumerate}

（使用§3的命题2和习题10 a)。）

\begin{enumerate}
\def\labelenumi{\arabic{enumi})}
\setcounter{enumi}{3}
\tightlist
\item
  令 V 为维数 \(n \geq 2\) 的向量空间，\(\mathfrak{g} = \mathfrak{sl}(V)\)，\(G = \mathbf{PGL}(V)\) 与 g 的一个自同构群相识别。g 中的 \(sl_{2}\) 三元组赋予 V 一个忠实的 \(\mathfrak{sl}(2, k)\)-模结构，反之，任何此类结构都来自某个 \(sl_{2}\) 三元组；\(sl_{2}\) 三元组是主的当且仅当相应的 \(\mathfrak{sl}(2, k)\)-模是单模；两个 \(sl_{2}\) 三元组 G-共轭当且仅当相应的 \(\mathfrak{sl}(2, k)\)-模同构。由此推出，\(sl_{2}\) 三元组的 G-共轭类与满足下列条件的族 \((m_{1}, m_{2}, \ldots)\)（由整数 \(\geq 0\) 组成）一一对应
\end{enumerate}

\[
m _ {1} + 2 m _ {2} + 3 m _ {3} + \dots = n \quad \text { and } \quad m _ {1} <   n.
\]

¶ 5) 假设 g 是半单的。令 a 为 g 的一个子代数，它在 g 中约化、与 g 同秩，且含有 g 的一个主 \(sl_{2}\) 三元组。证明 \(a = g\)。

¶6) 假设 g 是绝对单的，并以 h 表示其 Coxeter 数。令 x 为 g 的幂零元素。证明 \((\operatorname{ad} x)^{2h-1} = 0\)，且 \((\operatorname{ad} x)^{2h-2} \neq 0\) 当且仅当 x 是主的。（化归到 g 分裂且 x 含于命题10的子代数 \(n_{+}\) 中的情形。重复命题10的证明。）

¶7) 假设 g 是半单的。令 x 为 g 的幂零元素。则 x 是主的当且仅当 x 含于 g 的唯一的 Borel 子代数中。（化归到 k 代数闭的情形。使用命题10以及§3第3节的命题10。）

\begin{enumerate}
\def\labelenumi{\arabic{enumi})}
\setcounter{enumi}{7}
\item
  半单李代数具有 \(\mathfrak{sl}_2\) 三元组当且仅当它 \(\neq 0\) 且具有 Borel 子代数。
\item
  假设 g 是半单的。令 N（分别地 P）为 g 的幂零（分别地主幂零）元素的集合。
\end{enumerate}

\begin{enumerate}
\def\labelenumi{\alph{enumi})}
\item
  证明 \(\mathrm{P}\) 是 \(\mathbf{N}\) 在 Zariski 拓扑中的开子集（使用习题6）。
\item
  假设 g 可分裂。证明 P 在 N 中稠密。（使用命题10以及§10定理1的推论2。）
\end{enumerate}

\begin{enumerate}
\def\labelenumi{\arabic{enumi})}
\setcounter{enumi}{9}
\tightlist
\item
  假设 \(\mathfrak{g}\) 是可分裂半单的。令 \((x,h,y)\) 为一个 \(\mathfrak{sl}_2\) 三元组，位于 \(\mathfrak{g}\) 中。
\end{enumerate}

\begin{enumerate}
\def\labelenumi{\alph{enumi})}
\item
  证明存在一个分裂嘉当子代数 \(\mathfrak{h}\)，它含于 \(\mathfrak{g}\) 且包含 \(h\)。（使用§3的习题10 b)。）
\item
  取 \(\mathfrak{h}\) 如 \(a\) )。证明 \(h \in \mathfrak{h}_{\mathbf{Q}}\)，且存在 \(\mathrm{R}(\mathfrak{g}, \mathfrak{h})\) 的基 B，使得 \(\alpha(h) \in \{0, 1, 2\}\) 对所有 \(\alpha \in \mathrm{B}\) 成立（参见命题5）。元素 \(x\) 属于 \(\mathfrak{g}\) 的由诸 \(\mathfrak{g}^{\alpha}\)（\(\alpha \in \mathrm{B}\)）生成的子代数。
\item
  由 a)、b) 以及 Jacobson–Morozov 定理，给出“g 的每个幂零元素都含于某个 Borel 子代数”的一个新证明（参见§10定理1的推论2）。
\end{enumerate}

¶11) 令 \((x,h,y)\) 为一个主 \(\mathfrak{sl}_2\) 三元组，位于半单李代数 \(\mathfrak{g}\) 中。赋予 \(\mathfrak{g}\) 由该三元组定义的 \(\mathfrak{sl}(2,k)\)-模结构。证明如此定义的模同构于 \(\bigoplus_{i=1}^{l} V(2k_i - 2)\)，其中 \(k_i\) 是 \(\mathfrak{g}\) 上不变多项式函数代数的特征次数。（化归到 \(\mathfrak{g}\) 是可分裂单李代数的情形。使用§8第3节定理1的推论1，以及第VI章§4习题6 c。）\(^{20}\)

¶ 12) 假设 g 是半单的。令 \(x \in g\)，并令 s（分别地 n）为 x 的半单（分别地幂零）分量。令 \(a_{x}\)（分别地 \(a_{s}\)）为 x（分别地 s）在 g 中的交换子。

\begin{enumerate}
\def\labelenumi{\alph{enumi})}
\item
  证明 \(n\) 是半单代数 \(\mathcal{D}(\mathfrak{a}_s)\) 的幂零元素，且 \(n\) 在 \(\mathfrak{a}_s\) 中的交换子等于 \(\mathfrak{a}_x\)。由此推出 \(\dim \mathfrak{a}_x < \dim \mathfrak{a}_s\)（若 \(n \neq 0\)，即若 \(x\) 不是半单的）。
\item
  证明 \(\dim \mathfrak{a}_x = \mathrm{rk}(\mathfrak{g})\) 当且仅当 \(n\) 是 \(\mathcal{D}(\mathfrak{a}_s)\) 的主幂零元素。
\item
  置 \(\mathrm{G} = \mathrm{Aut}_{e}(\mathfrak{g})\)。证明对所有 \(\lambda \in k\)，存在 \(\sigma_{\lambda} \in G\) 使得 \(\sigma_{\lambda}x = s + \lambda^{2}n\)。（证明：若 \(n \neq 0\)，则存在一个 \(sl_{2}\) 三元组（在 \(a_{s}\) 中），其第一分量为 n，由此得到一个同态 \(\varphi : \mathbf{SL}(2, k) \to \mathbf{G}\)；取 \(\sigma_{\lambda}\) 为在 \(\varphi\) 下 \(\mathbf{SL}(2, k)\) 的适当对角元素的像。）由此推出 s 属于 G.x 在 Zariski 拓扑中的闭包。
\end{enumerate}

\(d\) ) 证明，若 \(x\) 不是半单的，则 \(x\) 不属于 G.s 在 Zariski 拓扑中的闭包（利用不等式 \(\dim \mathfrak{a}_x < \dim \mathfrak{a}_s\)，参见 \(a\) )）。

\begin{enumerate}
\def\labelenumi{\alph{enumi})}
\setcounter{enumi}{4}
\tightlist
\item
  假设 \(k\) 代数闭。证明下列性质等价：
\end{enumerate}

\begin{enumerate}
\def\labelenumi{(\roman{enumi})}
\item
  \(x\) 是半单的；
\item
  G.x 在 g 中是 Zariski 拓扑意义下闭的。
\end{enumerate}

（蕴含关系 (ii) \(\Longrightarrow\) (i) 由 c) 得到。若 (i) 成立且 \(x'\) 属于 G.x 的闭包，则§8的习题15表明 \(s'\)（\(x'\) 的半单分量）属于 G.x，于是 \(x'\) 属于 G. \(s'\) 的闭包；对 \(x'\) 与 \(s'\) 应用 d) 即得结论。）

\(f\) ) 假设 \(k\) 代数闭。令 \(\mathrm{F}_x\) 为由元素 \(y\in \mathfrak{g}\)（满足 \(f(x) = f(y)\)，对每个不变多项式函数 \(f\)（\(\mathfrak{g}\) 上的））组成的集合；我们有 \(y\in \mathrm{F}_x\) 当且仅当 \(y\) 的半单分量 G-共轭于 \(s\)（§8习题17）。证明 \(\mathrm{F}_x\) 是 G 的有限多条轨道的并，且轨道条数 \(\leq 3^{l(x)}\)，其中 \(l(x)\) 是 \(\mathcal{D}(\mathfrak{a}_s)\) 的秩。这些轨道中只有一条是闭的：\(s\) 的轨道；只有一条在 \(\mathrm{F}_x\) 中是开的：由元素 \(y\in \mathrm{F}_x\)（满足 \(\dim \mathfrak{a}_y = \operatorname {rk}(\mathfrak{g})\)）组成的那条。

\begin{enumerate}
\def\labelenumi{\arabic{enumi})}
\setcounter{enumi}{12}
\tightlist
\item
  假设 \(\mathfrak{g}\) 是半单的。
\end{enumerate}

\begin{enumerate}
\def\labelenumi{\alph{enumi})}
\item
  令 \((x, h, y)\) 为 g 中的主 \(sl_{2}\) 三元组，并令 b 为包含 x 的 Borel 子代数（习题7）。证明 b 含于 Im ad x。
\item
  假设 k 代数闭。证明对 g 中任意元素 z，存在 \(x, t \in g\)，其中 x 为主幂零元，使得 \(z = [x, t]\)（对包含 z 的一个 Borel 子代数应用 a)）。
\end{enumerate}

\begin{enumerate}
\def\labelenumi{\arabic{enumi})}
\setcounter{enumi}{13}
\tightlist
\item
  假设 \(\mathfrak{g}\) 是半单的。令 \(\mathfrak{p}\) 为 \(\mathfrak{g}\) 的抛物子代数，并令 \(f_{1}\)、\(f_{2}\) 为从 \(\mathfrak{g}\) 到某个有限维李代数的两个同态。证明 \(f_{1}|\mathfrak{p} = f_{2}|\mathfrak{p}\) 蕴含 \(f_{1} = f_{2}\)。（化归到 \(\mathfrak{g}\) 分裂的情形，再化归到 \(\mathfrak{g} = \mathfrak{sl}(2,k)\) 的情形，并使用第1节的引理1。）
\end{enumerate}

¶ 15) 假设 g 是半单的。

\begin{enumerate}
\def\labelenumi{\alph{enumi})}
\item
  令 \(x\) 为 \(\mathfrak{g}\) 的幂零元素。证明 \(x\) 含于 \(\operatorname{Im}(\operatorname{ad} x)^2\)（使用命题2）。由此推出 \((\operatorname{ad} x)^2 = 0\) 蕴含 \(x = 0\)。
\item
  令 \((x,h,y)\) 为一个 \(\mathfrak{sl}_2\) 三元组（在 \(\mathfrak{g}\) 中），并令 \(\mathfrak{s} = kx\oplus kh\oplus ky\)。证明下列条件等价：
\end{enumerate}

\begin{enumerate}
\def\labelenumi{(\roman{enumi})}
\item
  \(\operatorname{Im}(\operatorname{ad} x)^{2}=k.x;\)
\item
  \(\mathfrak{s}\)-模 \(\mathfrak{g} / \mathfrak{s}\) 是维数为 1 或 2 的单模之和；
\item
  \(\mathrm{ad}_{\mathfrak{g}}h\) 的异于 0、1 和 -1 的特征值只有 2 和 -2，且它们的重数都等于 1。
\end{enumerate}

\begin{enumerate}
\def\labelenumi{\alph{enumi})}
\setcounter{enumi}{2}
\item
  假设 g 是可分裂单李代数。令 h 为 g 的分裂嘉当子代数，B 为 \(\mathrm{R}(\mathfrak{g},\mathfrak{h})\) 的基，\(\gamma\) 为 \(\mathrm{R}(\mathfrak{g},\mathfrak{h})\) 相对于 B 的最高根。令 \((x,h,y)\) 为一个 \(sl_{2}\) 三元组，满足 \(h\in h_{Q}\) 且 \(\alpha(h)\geq0\) 对所有 \(\alpha\in B\)（参见命题5）。证明 b) 中的条件 (i)、(ii)、(iii) 成立当且仅当 \(h=H_{\gamma}\)，此时 \(x\in g^{\gamma}\) 且 \(y\in g^{-\gamma}\)。
\item
  保留 c) 的假设，并置 \(\mathrm{G} = \mathrm{Aut}_{0}(\mathfrak{g})\)。证明满足 (i)、(ii) 和 (iii) 的 \(sl_{2}\) 三元组是 G-共轭的（使用习题10）。证明：若 \(x\) 是 \(\mathfrak{g}\) 的满足 (i) 的非零幂零元素，则 \(G.x\) 在 Zariski 拓扑中的闭包等于 \(\{0\} \cup G.x\)。
\end{enumerate}

\begin{enumerate}
\def\labelenumi{\arabic{enumi})}
\setcounter{enumi}{15}
\tightlist
\item
  令 \((x,h,y)\) 为 g 中的一个 \(sl_{2}\) 三元组。证明：若 -2 是 k 中的平方，则元素 x-y 与 h 经 \(\operatorname{Aut}_{e}(\mathfrak{g})\) 的一个元素共轭（化归到 \(\mathfrak{g}=\mathfrak{sl}(2,k)\) 的情形）。由此推出：若 g 半单，则 x-y 是半单的，且它是正则的当且仅当 h 是正则的。
\end{enumerate}

¶ 17) 令 \((x, h, y)\) 为一个主 \(\mathfrak{sl}_2\) 三元组，位于半单代数 \(\mathfrak{g}\) 中。对所有 \(i \in \mathbf{Z}\)，以 \(\mathfrak{g}_i\) 表示 ad \(h\) 相应于特征值 \(i\) 的特征子空间；我们有 \(\mathfrak{g} = \bigoplus_{i \in \mathbf{Z}} \mathfrak{g}_i\)，且 \(\mathfrak{g}_i = 0\)（当 \(i\) 为奇数）。直和 \(\mathfrak{b}\)——诸 \(\mathfrak{g}_i\)（\(i \geq 0\)）之和——是 \(\mathfrak{g}\) 的 Borel 子代数，而 \(\mathfrak{g}_0\) 是嘉当子代数。

\begin{enumerate}
\def\labelenumi{\alph{enumi})}
\item
  证明对所有 \(z \in \mathfrak{b}\)，\(y + z\) 在 \(\mathfrak{g}\) 中的交换子的维数为 \(l = \operatorname{rk}(\mathfrak{g})\)。
\item
  令 I 为 \(\mathfrak{g}\) 上的不变多项式函数代数，\(\mathrm{P}_1, \ldots, \mathrm{P}_l\) 为生成 I 的 I 中齐次元素；置 \(\deg(\mathrm{P}_i) = k_i = m_i + 1\)。证明 \(\mathfrak{c}\)（\(x\) 在 \(\mathfrak{g}\) 中的交换子）具有基 \(x_1, \ldots, x_l\)，其中 \(x_i \in \mathfrak{g}_{2m_i}\)（参见习题11）。
\item
  令 \(i \in \{1, l\}\)，并令 \(J_i\)（分别地 \(K_i\)）为由 \(j \in (1, l)\)（满足 \(m_j = m_i\)（分别地 \(m_j < m_i\)））组成的集合。令 \(f_i \in k[X_1, \ldots, X_l]\) 为使得
\end{enumerate}

\[
f _ {i} \left(a _ {1}, \dots , a _ {l}\right) = \mathrm{P} _ {i} \left(y + \sum_ {j = 1} ^ {j = l} a _ {j} x _ {j}\right) \quad \text { for } \left(a _ {j}\right) \in k ^ {l}.
\]

证明 \(f_{i}\) 是线性型 \(\mathrm{L}_i\)（关于诸 \(\mathrm{X}_j\)，\(j \in \mathrm{J}_i\)）与一个多项式（关于诸 \(\mathrm{X}_j\)，\(j \in \mathrm{K}_i\)）之和。（若 \(t \in k^*\)，则 \(\mathfrak{g}\) 上等于 \(t^i\)（在 \(\mathfrak{g}_i\) 上）的自同构属于 \(\operatorname{Aut}_0(\mathfrak{g})\)，它把 \(y\) 变为 \(t^{-2}y\)，把 \(x_j\) 变为 \(t^{2m_j}x_j\)。利用 \(\mathrm{P}_i\) 在该自同构下的不变性。）

\(d)\) 令 \(\mathrm{P}\) 为从 \(\mathfrak{g}\) 到 \(k^l\) 的、由诸 \(\mathrm{P}_i\) 定义的映射。若 \(z \in \mathfrak{g}\)，以 \(\mathrm{D}_z\mathrm{P} : \mathfrak{g} \to k^l\) 表示 \(\mathrm{P}\) 在 \(z\) 处的切线性映射（第VII章附录I第2节）。

证明 \(\mathfrak{c}\cap\operatorname{Im}\operatorname{ad}(y-x)=0\)（相对于由所给 \(sl_{2}\) 三元组生成的子代数，把 g 分解为单子模的直和）。由此推出 \(D_{y-x}P\) 在 c 上的限制是从 c 到 \(k^{l}\) 的同构（利用前一习题及§8的习题18 a)）。利用这一结果证明，c) 中定义的线性型 \(L_{i}\) 的行列式 \(\neq0\)，并由此推出下列结果：

\(d_{1})\) 多项式 \(f_{1},\ldots ,f_{l}\) 代数无关，且生成 \(k[\mathrm{X}_1,\dots ,\mathrm{X}_l]\)；

\(d_{2})\) 映射 \(z \mapsto \mathrm{P}(y + z)\)（从 \(\mathfrak{c}\) 到 \(k^{l}\)）是双射多项式映射，且其逆映射也是多项式映射；

\(d_{3})\) 对所有 \(z\in y + \mathfrak{c}\)，线性映射 \(\mathrm{D}_z\mathrm{P}|\mathfrak{c}\) 的秩为 \(l\)。

特别地，映射 \(\mathrm{P}:\mathfrak{g}\to k^l\) 是满射。

\begin{enumerate}
\def\labelenumi{\alph{enumi})}
\setcounter{enumi}{4}
\tightlist
\item
  若 \(k\) 代数闭，则 \(\mathfrak{g}\) 中每个交换子维数为 \(l\) 的元素都经 \(\mathrm{Aut}_0(\mathfrak{g})\) 共轭于 \(y + \mathfrak{c}\) 的唯一元素。
\end{enumerate}

\(f\) ) 给出一个没有主 \(\mathfrak{sl}_2\) 三元组且映射 \(\mathrm{P}\) 非满射的单李代数的例子（取 \(k = \mathbf{R}\) 及 \(l = 1\)）。

